% Section-only TeX block for the MAH paper.
% Suggested bibkeys below. Replace them if your .bib uses different keys.
%
% Hopenhayn1992
% EricsonPakes1995
% KletteKortum2004
% AkcigitKerr2018
% AroraFosfuriGambardella2001
% AroraMerges2004
% GansStern2000
% GansStern2003
% Lakdawalla2018
% AcemogluLinn2004
% JiaMaYangZhang2023
% FonsRosenRoldanBlancoSchmitz2024

\section{Introduction}

A recurring theme in the economics of innovation is that invention and commercialization need not be organized inside the same firm. This point is especially salient in the pharmaceutical industry, where turning a scientifically promising idea into a marketable drug requires not only discovery, but also regulatory execution, quality accountability, and reliable production. In such an environment, institutional design matters not merely because it changes prices or demand, but because it changes which organizational routes from invention to commercialization are feasible, enforceable, and scalable.

This paper studies China's Marketing Authorization Holder (MAH) reform through that lens. Our question is not whether MAH is a generic pro-innovation policy, nor whether it directly raises scientific productivity. Instead, we ask whether MAH changes the organizational route through which research-originated original-drug ideas are brought to market, and thereby increases the number of ideas that can be realized as marketable products. That is the paper's central object.

That question is well suited to the Chinese pharmaceutical industry. The industry has long been shaped by a structural tension between a growing upstream research margin and a historically tighter linkage between authorization and internal production. Under the pre-reform regime, a research-originating firm without internal manufacturing capacity faced substantial difficulty in bringing an original-drug idea to market while retaining project control and authorization status. In practice, a scientifically promising project could fail not because it lacked scientific merit, but because commercialization required either downstream integration or a transfer of effective control to an integrated producer. MAH changes that institutional environment by explicitly allowing the holder to retain authorization while commissioning production from a qualified external manufacturer under a legally recognized structure with clear quality responsibility and standardized compliance requirements.

The economic implication follows directly. MAH lowers the governance and compliance friction of one specific organizational route: external commercialization by a research-originating firm that retains authorization status. We denote that route-specific friction by $\tau_{ext}$. This object is narrower and more informative than a generic policy dummy. It is not a stand-in for all regulatory costs in the industry, nor for demand, scientific productivity, or general business conditions. It captures the extra organizational and compliance burden of implementing an original-drug idea through a qualified external manufacturer while the innovator keeps the project and holder status. If MAH lowers $\tau_{ext}$, then the expected net value of externally commercialized original-drug ideas rises, and so does the set of research-oriented firms for which commercialization becomes feasible.

This change in commercialization assignment is the main mechanism of the paper. Firm-boundary reorganization remains important, but it is no longer the only or even the first margin. A research-specialized firm may remain research-specialized and externalize commercialization to a manufacturing specialist. It may instead switch into an integrated form and internalize implementation. An integrated incumbent may continue to invent and implement internally. A project may also be transferred or abandoned. These are not separate stories. They are alternative organizational assignments for the same underlying problem: who ultimately commercializes a research-originated original-drug idea, and under which institutional conditions.

To study that problem, we build a dynamic industry model with three organizational forms. Type-$A$ firms are integrated: they conduct original-drug research and organize production internally. Type-$B$ firms are research-specialized: they generate original-drug ideas but do not operate as full internal manufacturers. Type-$C$ firms are manufacturing-specialized: they provide qualified production capability and serve the contract-manufacturing market. The model is embedded in a Hopenhayn-style stationary industry environment with endogenous entry, exit, and organizational switching. Original-drug innovation is modeled as a two-stage process: firms first generate ideas, and those ideas are then assigned to commercialization routes. MAH operates primarily on the second stage \cite{Hopenhayn1992,EricsonPakes1995,KletteKortum2004,AkcigitKerr2018,FonsRosenRoldanBlancoSchmitz2024}.

This modeling choice also matches the paper's identification boundary. The reform does not need to be interpreted as a direct increase in scientific productivity. It is sufficient that a previously difficult route---research-originated commercialization through qualified external manufacturing while retaining authorization---becomes more feasible and less costly. With standard firm-level and product-level outcomes, one can identify whether MAH raises the number of new original-drug products, the share of original drugs in realized products, the entry of research-oriented firms, or the use of entrusted production. One cannot directly observe a primitive governance friction from such outcomes alone. The role of $\tau_{ext}$ is therefore interpretive but disciplined, and it must be supported jointly by commercialization proxies, entrant composition, organizational-transition evidence, and producer-side outcomes.

The paper contributes to four literatures. First, it joins the dynamic heterogeneous-firm literature on entry, exit, and industry composition, following \citet{Hopenhayn1992} and \citet{EricsonPakes1995}, while adapting that industry shell to an organizational-boundary question rather than a pure productivity-selection question. It also draws on innovation-based firm dynamics in \citet{KletteKortum2004} and \citet{AkcigitKerr2018}, where innovative effort and firm composition jointly shape industry outcomes. Second, it connects to the literature on the organization of innovation and commercialization. In particular, it builds on the idea that innovators need not internalize all downstream production in order to create value \citep{AroraFosfuriGambardella2001,AroraMerges2004,GansStern2000,GansStern2003}. Third, it relates to the pharmaceutical innovation literature showing that policy often affects innovation through expected returns, delay, horizon, coverage, and regulatory design rather than through factory-level productivity alone \citep{AcemogluLinn2004,Lakdawalla2018}. Fourth, it speaks directly to the Chinese pharmaceutical reform context, where \citet{JiaMaYangZhang2023} show that regulatory improvements can reshape firm composition and innovation outcomes. Finally, the paper borrows the distinction between invention and implementation from \citet{FonsRosenRoldanBlancoSchmitz2024}, while replacing startup acquisition with MAH-enabled commercialization through qualified external manufacturing.

Our contribution relative to those literatures is to show that, in a latecomer pharmaceutical industry, innovation policy may work not primarily by making research more productive in the scientific sense, but by changing the organizational routes through which ideas become products. In that sense, MAH is best understood not as a generic innovation wedge, but as a reform that opens and cheapens a specific commercialization route. That reframing also changes the empirical priority of the paper: the first-layer objects are commercialization assignment, conversion from research-originated ideas to realized products, the number of new original-drug products, and the share of original drugs in realized products; organizational transition matrices remain important, but they serve as complementary validation rather than as the sole organizing object.

The rest of the paper proceeds as follows. Section 2 describes the institutional setting and identification boundary. Section 3 develops the model. Section 4 characterizes equilibrium and states the main comparative-static predictions. Section 5 presents the empirical design and estimation strategy. Section 6 reports the results, counterfactuals, and welfare implications. Section 7 concludes.
